\begin{frame}
\frametitle{Moduly s dynamickou pamětí}

\centering

\includegraphics[width=0.85\linewidth]{generated/memory/dram/sdram-dimm.pdf}

\end{frame}

\begin{frame}
\frametitle{Paměť v dnešních osobních počítačích}

{
\centering

\includegraphics[width=0.85\linewidth]{generated/memory/pc-memory-to-numa.pdf}

}

\vskip 2mm

MC - Memory controller -- paměťový kontrolér zajišťuje řízení čtení a zápisů do SDRAM. Dále obstarává obnovu informace (refreshing) každé paměťové buňky jednou za 64 ms.

\end{frame}


\begin{frame}
\frametitle{Nejčastěji používané typy dynamických pamětí}

\begin{itemize}
\item \textbf{SDRAM} -- hodinová frekvence až 100 MHz, 2.5V, synchronní přenos dat na hodinové hraně
\item \textbf{DDR SDRAM} -- přenos dat na obou hodinových hranách, 2.5V, V/V hodiny až 100-200 MHz, 0.2-0.4 GT/s (miliard přenosů za skundu)
\item \textbf{DDR2 SDRAM} -- menší spotřeba, 1.8V, frekvence až 400 MHz, 0.8 GT/s
\item \textbf{DDR3 SDRAM} -- ještě nižsí spotřeba při 1.5V, frekvence až 800 MHz, 1.6 GT/s
\item \textbf{DDR4 SDRAM} -- 1.05 -- 1.2V,  V/V svěrnicové hodiny 1.2 GHz, 2.4 GT/s
\item \textbf{DDR5 SDRAM} -- 1.1V, až 6.4 GT/s
\end{itemize}

Všechny tyto druhy jsou převážně optimalizované na propustnost, nikoliv náhodný přístup, latence 20 až 35 ns.

\end{frame}
